\documentclass[10pt,a4paper]{amsart}
\usepackage{amsmath,amssymb,amsthm,graphicx}
% Retained preamble for the Stacks Project chapter snapshot in this directory.
% The chapter's own preamble.tex is a copy of the shared Stacks Project preamble; the block
% below is sliced MECHANICALLY, line for line, from it. Nothing was added, removed or
% redefined; the wrapper class, the xr-hyper/hyperref wiring and the externaldocument list
% are omitted because the assembler supplies the document class.
%
% lines taken from preamble.tex: 3, 5, 9, 11, 12, 156, 157, 158, ... (51 lines total)
%
% package coverage restored from the chapter's own preamble (2026-10-13): xy, xy-2cell

\usepackage[all]{xy}
\xyoption{2cell}
\UseAllTwocells

\newenvironment{reference}{\comment}{\endcomment}
\newenvironment{slogan}{\comment}{\endcomment}
\newenvironment{history}{\comment}{\endcomment}
\theoremstyle{plain}
\newtheorem{theorem}[subsection]{Theorem}
\newtheorem{proposition}[subsection]{Proposition}
\newtheorem{lemma}[subsection]{Lemma}
\theoremstyle{definition}
\newtheorem{definition}[subsection]{Definition}
\newtheorem{example}[subsection]{Example}
\newtheorem{exercise}[subsection]{Exercise}
\newtheorem{situation}[subsection]{Situation}
\theoremstyle{remark}
\newtheorem{remark}[subsection]{Remark}
\newtheorem{remarks}[subsection]{Remarks}
\numberwithin{equation}{subsection}
\def\lim{\mathop{\mathrm{lim}}\nolimits}
\def\colim{\mathop{\mathrm{colim}}\nolimits}
\def\Spec{\mathop{\mathrm{Spec}}}
\def\Hom{\mathop{\mathrm{Hom}}\nolimits}
\def\Ext{\mathop{\mathrm{Ext}}\nolimits}
\def\SheafHom{\mathop{\mathcal{H}\!\mathit{om}}\nolimits}
\def\SheafExt{\mathop{\mathcal{E}\!\mathit{xt}}\nolimits}
\def\Sch{\mathit{Sch}}
\def\Mor{\mathop{\mathrm{Mor}}\nolimits}
\def\Ob{\mathop{\mathrm{Ob}}\nolimits}
\def\Sh{\mathop{\mathit{Sh}}\nolimits}
\def\NL{\mathop{N\!L}\nolimits}
\def\CH{\mathop{\mathrm{CH}}\nolimits}
\def\proetale{{pro\text{-}\acute{e}tale}}
\def\etale{{\acute{e}tale}}
\def\QCoh{\mathit{QCoh}}
\def\Ker{\mathop{\mathrm{Ker}}}
\def\Im{\mathop{\mathrm{Im}}}
\def\Coker{\mathop{\mathrm{Coker}}}
\def\Coim{\mathop{\mathrm{Coim}}}
\DeclareMathSymbol{\boxtimes}{\mathbin}{AMSa}{"02}
\def\QCohstack{\mathcal{QC}\!\mathit{oh}}
\def\Cohstack{\mathcal{C}\!\mathit{oh}}
\def\Spacesstack{\mathcal{S}\!\mathit{paces}}
\def\Quotfunctor{\mathrm{Quot}}
\def\Hilbfunctor{\mathrm{Hilb}}
\def\Curvesstack{\mathcal{C}\!\mathit{urves}}
\def\Polarizedstack{\mathcal{P}\!\mathit{olarized}}
\def\Complexesstack{\mathcal{C}\!\mathit{omplexes}}
\def\Pic{\mathop{\mathrm{Pic}}\nolimits}
\def\Picardstack{\mathcal{P}\!\mathit{ic}}
\def\Picardfunctor{\mathrm{Pic}}
\def\Deformationcategory{\mathcal{D}\!\mathit{ef}}
\usepackage{cleveref}
\crefname{originallemma}{original lemma}{original lemmas}
\Crefname{originallemma}{Original lemma}{Original lemmas}
\makeatletter
\newlabel{algebra-lemma-universal-omega}{{00RO}{source}{}{sourcepointer-1}{}}
\newlabel{algebra-lemma-universal-omega@cref}{{[originallemma][0][]00RO}{[1][0][]source}{}{}{}}
\makeatother
\crefname{originallemma}{original lemma}{original lemmas}
\Crefname{originallemma}{Original lemma}{Original lemmas}
\makeatletter
\newlabel{defos-lemma-choices}{{08S7}{source}{}{sourcepointer-2}{}}
\newlabel{defos-lemma-choices@cref}{{[originallemma][0][]08S7}{[1][0][]source}{}{}{}}
\makeatother
\crefname{originallemma}{original lemma}{original lemmas}
\Crefname{originallemma}{Original lemma}{Original lemmas}
\makeatletter
\newlabel{defos-lemma-deform-module}{{08LI}{source}{}{sourcepointer-3}{}}
\newlabel{defos-lemma-deform-module@cref}{{[originallemma][0][]08LI}{[1][0][]source}{}{}{}}
\makeatother
\crefname{originallemma}{original lemma}{original lemmas}
\Crefname{originallemma}{Original lemma}{Original lemmas}
\makeatletter
\newlabel{defos-lemma-huge-diagram}{{08S5}{source}{}{sourcepointer-4}{}}
\newlabel{defos-lemma-huge-diagram@cref}{{[originallemma][0][]08S5}{[1][0][]source}{}{}{}}
\makeatother
\crefname{originallemma}{original lemma}{original lemmas}
\Crefname{originallemma}{Original lemma}{Original lemmas}
\makeatletter
\newlabel{example-rings}{{0DY1}{source}{}{sourcepointer-5}{}}
\newlabel{example-rings@cref}{{[originallemma][0][]0DY1}{[1][0][]source}{}{}{}}
\makeatother
\crefname{originallemma}{original lemma}{original lemmas}
\Crefname{originallemma}{Original lemma}{Original lemmas}
\makeatletter
\newlabel{more-morphisms-lemma-deform}{{063Y}{source}{}{sourcepointer-6}{}}
\newlabel{more-morphisms-lemma-deform@cref}{{[originallemma][0][]063Y}{[1][0][]source}{}{}{}}
\makeatother
\title{Stacks, Examples of Deformations: lemma-rings-TI}
\author{The Stacks Project Authors}
\date{stacks-project:examples-defos.tex:local-snapshot-20261012; GNU FDL 1.2 or later}
\begin{document}
\maketitle
\noindent\textit{Editorial note.} Statement span is lines 1071-1079 of sources/stacks-examples-defos/pinned-original-source.tex; proof span is lines 1081-1107, the complete author proof (25 body lines).\par\medskip
\noindent\textit{Editorial note.} Navigation stubs: 6 cited-but-not-included labels are supplied as non-load-bearing pointers with their pinned tags where the snapshot tag index has them.\par\medskip
\noindent\textit{Editorial note.} extra\_packages is empty on purpose: the retained preamble already loads xy with options and a second declaration clashes in the assembler wrapper.\par\medskip
\noindent\textit{Editorial note.} Text is retained verbatim, including the author's own typos and placeholders.\par\medskip
\noindent\textit{Editorial note.} Queue index 22 of queue\_examples-defos.json (pb=1146, stmt 1071-1079, proof 1081-1107).\par\medskip
\noindent\textit{Allowed standard prerequisites (editorial).}\begin{itemize}
\item Explicit prerequisite boundary. The cited inputs are prerequisites; the argument in the retained span is not replaced by any cited result.
\end{itemize}\medskip
\ifdefined\hypertarget\hypertarget{sourcepointer-1}{}\fi
\noindent\textit{Source reference.} Original lemma 00RO: sources/stacks-examples-defos/pinned-original-source.tex, source lines 1107--1107. Cited by the retained statement or proof but not part of the collected spans; supplied only as a navigation pointer. The argument does not depend on its content.\par\medskip
\ifdefined\hypertarget\hypertarget{sourcepointer-2}{}\fi
\noindent\textit{Source reference.} Original lemma 08S7: sources/stacks-examples-defos/pinned-original-source.tex, source lines 1107--1107. Cited by the retained statement or proof but not part of the collected spans; supplied only as a navigation pointer. The argument does not depend on its content.\par\medskip
\ifdefined\hypertarget\hypertarget{sourcepointer-3}{}\fi
\noindent\textit{Source reference.} Original lemma 08LI: sources/stacks-examples-defos/pinned-original-source.tex, source lines 1107--1107. Cited by the retained statement or proof but not part of the collected spans; supplied only as a navigation pointer. The argument does not depend on its content.\par\medskip
\ifdefined\hypertarget\hypertarget{sourcepointer-4}{}\fi
\noindent\textit{Source reference.} Original lemma 08S5: sources/stacks-examples-defos/pinned-original-source.tex, source lines 1107--1107. Cited by the retained statement or proof but not part of the collected spans; supplied only as a navigation pointer. The argument does not depend on its content.\par\medskip
\ifdefined\hypertarget\hypertarget{sourcepointer-5}{}\fi
\noindent\textit{Source reference.} Original lemma 0DY1: sources/stacks-examples-defos/pinned-original-source.tex, source lines 1031--1040. Cited by the retained statement or proof but not part of the collected spans; supplied only as a navigation pointer. The argument does not depend on its content.\par\medskip
\ifdefined\hypertarget\hypertarget{sourcepointer-6}{}\fi
\noindent\textit{Source reference.} Original lemma 063Y: sources/stacks-examples-defos/pinned-original-source.tex, source lines 1107--1107. Cited by the retained statement or proof but not part of the collected spans; supplied only as a navigation pointer. The argument does not depend on its content.\par\medskip
\begin{lemma}
\label{lemma-rings-TI}
In Example \ref{example-rings} let $P$ be a $k$-algebra. Then
$$
T\Deformationcategory_P = \text{Ext}^1_P(\NL_{P/k}, P)
\quad\text{and}\quad
\text{Inf}(\Deformationcategory_P) = \text{Der}_k(P, P)
$$
\end{lemma}

\begin{proof}
Recall that $\text{Inf}(\Deformationcategory_P)$ is the set of
automorphisms of the trivial deformation
$P[\epsilon] = P \otimes_k k[\epsilon]$ of $P$ to $k[\epsilon]$
equal to the identity modulo $\epsilon$.
By Deformation Theory, Lemma \ref{defos-lemma-huge-diagram}
this is equal to $\Hom_P(\Omega_{P/k}, P)$ which in turn is
equal to $\text{Der}_k(P, P)$ by
Algebra, Lemma \ref{algebra-lemma-universal-omega}.

\medskip\noindent
Recall that $T\Deformationcategory_P$ is the set of isomorphism classes
of flat deformations $Q$ of $P$ to $k[\epsilon]$, more precisely,
the set of isomorphism classes of $\Deformationcategory_P(k[\epsilon])$.
Recall that a $k[\epsilon]$-algebra $Q$ with $Q/\epsilon Q = P$
is flat over $k[\epsilon]$ if and only if
$$
0 \to P \xrightarrow{\epsilon} Q \to P \to 0
$$
is exact. This is proven in More on Morphisms, Lemma
\ref{more-morphisms-lemma-deform} and more generally in
Deformation Theory, Lemma \ref{defos-lemma-deform-module}.
Thus we may apply
Deformation Theory, Lemma \ref{defos-lemma-choices}
to see that the set of isomorphism classes of such
deformations is equal to $\text{Ext}^1_P(\NL_{P/k}, P)$.
\end{proof}

\end{document}
